\section{两条轴：离散化与容差}

\noindent\textbf{入口}：库里 97 条数学都活在两条轴上——一条是\textbf{连续 $\to$ 离散}（定理变成有限步算法），
一条是\textbf{精确 $\to$ 容差}（算法落到浮点）。说「这个算法对」之前，必须先说清它在这两条轴上各掉了多少。
本册的每条判据，都是这两条轴上的一个坐标。

\subsection{轴一：离散——定理怎么变成有限步}

连续对象换成有限样本、网格或基函数之后，误差按步长走：

\[
\norm{u-u_h}\le C\,h^{p},
\]

$h$ 是网格尺寸，$p$ 是方法阶。这条不等式是全册最重要的形状：\textbf{它把「对不对」换成了「按什么速率收敛」，后者才是可测的。}

\begin{center}
\begin{tabular}{L{3.4cm}L{5.2cm}L{6.2cm}}
\toprule
离散手段 & 判据 & 失效点 \\
\midrule
网格 + 有限元 & $h$ 减半，误差降 $2^{p}$，$p=$ 基函数多项式次数 & $h$ 未进渐近区时测出的「阶」不是阶 \\
采样 & 采样率 $>2f_{\max}$；折叠频率 $f \bmod f_s$ & 欠采样后混叠不可逆，无法从样本恢复 \\
求积 & $n$ 点 Gauss 对 $2n-1$ 次多项式精确 & 被积函数有间断，阶掉到 1 \\
截断（DFT / 级数） & 残差按 $\abs{\lambda_N}$ 或系数衰减 & 系数的衰减由光滑度决定，不由谱决定 \\
\bottomrule
\end{tabular}
\end{center}

\noindent\textbf{一条经验}：阶的测量必须用\textbf{三个点}（$h$、$h/2$、$h/4$），比值落在 $[2^{p}-\delta,\ 2^{p}+\delta]$ 才算测到；
单点比值乱跳只说明还没进渐近区。

\subsection{轴二：容差——「相等」必须换成「$\le \eps$」}

浮点没有相等，只有相对与绝对容差。所以本册的判据一律写成不等式，且必须带三要素：
\textbf{量、阈值、阈值的量级依据}。缺一条就不写。

\begin{center}
\begin{tabular}{L{4.6cm}L{9.8cm}}
\toprule
判据形状 & 用法与陷阱 \\
\midrule
$\norm{x-y}\le \eps_{\text{mach}}\norm{x}$ & 只对同量级有效；量级差三个数量级时相对误差已经不适用 \\
$\norm{x-y}\le a\,\eps_{\text{mach}}+r\,\norm{y}$ & 绝对 + 相对双阈值，库里比较函数统一用这个形状 \\
$\norm{A^{H}(b-Ax)}\le tol\cdot\norm{A^{H}b}$ & 最小二乘判\textbf{后向}残差；前向误差还要除以条件数 \\
$\kappa(A)\,\eps_{\text{mach}}\ll 1$ & 条件数放大：算法再好也救不了病态问题 \\
累积 $\sqrt{n}\,\eps$ 或 $n\,\eps$ & 随机游走型与最坏情形型，二者差一个 $\sqrt n$ \\
\bottomrule
\end{tabular}
\end{center}

\noindent\textbf{两条曲线}：$h\to0$ 时离散误差按 $h^{p}$ 降，但舍入被 $\kappa=O(h^{-2})$ 放大，
总误差在双对数图上出现\textbf{谷底}。\textbf{收敛阶测试的三个点必须都在谷底左侧}，否则测出来的不是 $p$。

\subsection{两条轴的交叉：为什么「数学正确」不等于「程序对」}

\begin{center}
\begin{tabular}{L{2.8cm}L{3.0cm}L{4.6cm}L{3.8cm}}
\toprule
定理 & 离散实现 & 浮点判据 & 失效点 \\
\midrule
正交归一 & MGS & $\norm{Q^{H}Q-I}\le 8n\eps$ & 不做重正交时正交性随 $n$ 退化 \\
可逆 & LU + 部分主元 & 残差小 & 残差小但解错：条件数大 \\
精确求积 & Simpson & 三次多项式精确 & 间断处阶降到 1 \\
谱分解 & QR 迭代 & 残差 $\le tol$ & 重根时特征向量不稳 \\
\bottomrule
\end{tabular}
\end{center}

\noindent\textbf{这张表的读法}：左边是数学，中间是实现，右边是判据，最后一列是判据失效的地方。
\textbf{四条缺一条就不算验证过}——只给左边的，是教材；只给中间的，是代码；四列齐全的，才是本册的条目。

\subsection{与「入口法」的关系（只借用两格）}

根目录 \texttt{LEARNING.md} 的六段在本册只用两格：
\textbf{入口}（这条数学在代码里长成什么形状）与\textbf{失效点}（状态词与数值极限）。
另外四格里的「生成关系」在本册有确定含义：\textbf{连续 $\to$ 离散 $\to$ 浮点}，每换一步就掉一次精度，
而每一掉都在上面两张表里有一行对应。这就是本册不写推导、只写判据的理由。
